\section{Abstract}

This thesis studies finite-horizon price tracking in a synthetic, zero-fee constant-product automated market maker (AMM) after a shock to an exogenous centralized-exchange (CEX) reference. It distinguishes exact swap mechanics from outcomes shaped by arbitrage capacity, liquidity depth, agent rules and event order. In an isolated pool, capacity $c$ can correct out-of-band log error by at most $2\log(1+c)$ per tick; for a 0.20 log shock and a 0.003 band, ideal correction takes 20 ticks at $c=0.005$, four at $c=0.03$ and two at $c=0.10$. The simulator includes noise, trend/panic, arbitrage and threshold-withdrawal liquidity-provider rules. It evaluates 13,500 main paths (300 seeds), 40,500 sensitivity paths (300 new seeds), and a 9,000-path post-hoc replay using archived seeds.

High-capacity arbitrage reduces mean cumulative excess-error area from 2.36 without arbitrage to 0.10, although its direction follows from the rule. At $c=0.03$, two-tick sustained return begins by the isolated-pool deadline in 49.0--84.0\,\% of paths at depth $\kappa=0.5$, 57.3--83.3\,\% at $\kappa=1$, and 98.3--100\,\% at $\kappa=2$; all recover by tick 20, and no main-run order reaches its capital limit. The E3 one-shot withdrawal trigger fires in every tested path, usually at tick 1. Supplementary maximum-response contrasts are positive only at shallow depth with fixed orders, disappear when orders scale with reserves, and shrink without panic. The replay finds that shallow static E1-high later departures (90/300 and 53/300 for shocks $-0.05$ and $-0.10$) are single panic-associated episodes, with none in matched no-panic paths. In E3, 98.2--99.6\,\% of recovered paths depart again, often repeatedly, with or without panic. Fixing arbitrage capacity at initial reserves reduces the shallow post-recovery excess-area contrast from 0.122--0.124 to 0.0036--0.0097. This post-hoc control covers one depth and conditions on recovery, so it does not establish a general causal effect.

The two-state map improves on persistence in some held-out conditions but does not outperform scalar AR(1), so its second state adds no demonstrated predictive value. A band-aware hybrid map reduces recursive replay error by about 1.7\,\%, but uses the realized future CEX path and is not a real-time forecast. The contribution is a reproducible, mechanism-focused assessment of the analytical benchmark and its limits. The results are conditional on a synthetic, zero-fee system and have no market-data validation; they establish neither market effects nor validated operating requirements. Synchronized event replay should record band entry, sustained return, later departures, orders, active reserves and available arbitrage capacity. These are candidate observables, not validated requirements; external assessment also needs realistic fees, gas, delays, trading costs, capital limits and arbitrage competition.

\textbf{Keywords:} automated market maker; arbitrage; price tracking; liquidity withdrawal; agent-based simulation; Monte Carlo
